% Autogenerated translation of LANGUAGE_MIXED_SYNTAX.md by Texpad
% To stop this file being overwritten during the typeset process, please move or remove this header

\documentclass[12pt]{book}
\usepackage{graphicx}
\usepackage{fontspec}
\usepackage[utf8]{inputenc}
\usepackage[a4paper,left=.5in,right=.5in,top=.3in,bottom=0.3in]{geometry}
\setlength\parindent{0pt}
\setlength{\parskip}{\baselineskip}
\setmainfont{Helvetica Neue}
\usepackage{hyperref}
\pagestyle{plain}
\begin{document}

\chapter*{Contract — Mixed OOP + Functional Syntax Design}

\section*{Goal}

Allow Contract source files to mix C\#-like OOP declarations (Contracts, \texttt{fn} methods, fields) and F\#-style functional expressions and declarations in the same file. Provide clear rules for parsing and semantics so both styles interoperate naturally.

\section*{Overview \& Motivation}

\begin{itemize}
\item Contract is primarily OOP/C\#-like: \texttt{Contract} types, fields, methods (\texttt{fn}), dot-qualified calls (e.g., \texttt{IO.println}), and \texttt{var} declarations.
\item You want to also allow functional-style code (lightweight expression syntax, piping, pattern matching, let-bindings) in the same file, enabling succinct functional idioms where appropriate.
\item Key constraint: a single parser (or coordinated parsers) must accept both syntaxes in context-sensitive ways without ambiguity. Diagnostics and tool support must remain accurate.
\end{itemize}

\section*{Design principles}

\begin{enumerate}
\item Contextual parsing: the syntax recognized depends on the surrounding construct. Inside a \texttt{fn} body or top-level module body, both expression syntaxes are allowed. Declarations (contracts, fields, fn headers) follow OOP syntax.
\item Minimal new syntax: reuse existing token set where possible. Add only a small set of functional sugar: \texttt{let} bindings, \texttt{$\vert$$>$} pipe operator, \texttt{match}/\texttt{with} pattern matching, lambda syntax \texttt{fun x -$>$ expr}, and optional infix operators for composition.
\item Deterministic parsing: ensure precedence \& associativity are well-defined and parsing decisions are local (no global backtracking across large ranges).
\item Interop semantics: functional constructs compile to the same IL/VM model — lambdas become synthetic functions/closures capturing locals, \texttt{let} maps to local vars, \texttt{match} compiles to switch-like bytecode.
\end{enumerate}

\section*{Proposed syntax additions}

\begin{itemize}
\item let binding (expression or statement position):

\begin{itemize}
\item \texttt{let x = expr} (creates a local named \texttt{x}), optionally \texttt{let rec} for recursion.
\item \texttt{let x: Int = expr} allowed
\end{itemize}
\item lambda: \texttt{fun x -$>$ expr} (single-arg lambdas), \texttt{fun (x,y) -$>$ expr} for tuples.
\item pipe operator: \texttt{expr $\vert$$>$ fn} (rewrites to \texttt{fn(expr)} during parsing/IR lowering)
\item match expression (expression position):

\begin{itemize}
\item \texttt{match expr with $\vert$ pattern -$>$ expr $\vert$ pattern -$>$ expr ...} where patterns are integer literals, identifiers, or \texttt{\_} wildcard.
\end{itemize}
\item sequence expressions: \texttt{expr; expr} is already present; \texttt{let} and \texttt{match} are expression-first where useful.
\end{itemize}

\section*{Parsing strategy}

Option A (recommended for speed):
- Implement a single grammar (nearley/PEG) that includes both OOP and functional constructs with contextual precedence rules (e.g., \texttt{let} recognized in expression/statement positions).
- Use the parser in the LSP for editor features and in the C compiler later (optionally by generating a C parser from the grammar or porting the grammar logic to the C parser module).

Option B (native C parser):
- Extend recursive-descent parser: expand \texttt{parse\_expr} to recognize \texttt{let}, \texttt{fun}, \texttt{match}, \texttt{$\vert$$>$} and \texttt{$\vert$} alternatives; treat these constructs only inside function bodies or top-level expression areas.
- Ensure \texttt{parse\_primary} handles \texttt{fun} and \texttt{match} and \texttt{let} as either expressions or statements depending on position.

\section*{Ambiguities and resolution}

\begin{itemize}
\item \texttt{let} vs \texttt{var}: both declare locals; adopt rule: \texttt{var} remains statement-level mutable declaration; \texttt{let} is expression-level and immutable by default (\texttt{let mut} can be added later).
\item \texttt{fn (...)} vs \texttt{fun}: \texttt{fn} is function declaration; \texttt{fun} is lambda. They are distinct tokens.
\item Pipe operator precedence: \texttt{$\vert$$>$} is low-precedence, right-associative.
\item Pattern matching vs switch: \texttt{match} is expression-first and can destructure; \texttt{switch} remains statement-level for legacy code.
\end{itemize}

\section*{Runtime/IL changes}

\begin{itemize}
\item Lambdas: represent as closures capturing environment. Implementation choices:

\begin{itemize}
\item Simple approach: compile \texttt{fun} to a method with synthetic name and capture variables as additional hidden parameters (no heap-allocated closure object). The caller/owner must pass captured values when calling the closure.
\item Future: add first-class closure objects on the heap. For v1, the synthetic-method approach keeps VM simple.
\end{itemize}
\item \texttt{let} -$>$ local variable allocation (similar to \texttt{var}). Map \texttt{let} to local\_map entries during emission.
\item \texttt{match} -$>$ lowered to equality checks and jump chains (like switch) or to a sequence of comparisons.
\end{itemize}

\section*{Examples}

Functional style inside fn body:

\texttt{
fn main() \{
  let inc = fun x -$>$ x + 1;
  let nums = [1,2,3];
  let doubled = nums $\vert$$>$ map (fun x -$>$ x * 2);
  IO.println("done");
\}
}

Pattern matching example:
\texttt{
let sign x = match x with
  $\vert$ 0 -$>$ 0
  $\vert$ n when n $>$ 0 -$>$ 1
  $\vert$ \_ -$>$ -1
}

Interop with OOP members:
\texttt{
Contract Program \{
  fn main() \{
    let msg = "Hello";
    IO.println(msg $\vert$$>$ String.toUpper);
  \}
\}
}

\end{document}
